\documentclass{article}
\usepackage{graphicx} % Required for inserting images
\usepackage{hyperref}
\usepackage{amsmath}
\usepackage{amssymb}
\hypersetup{
    colorlinks=true,
    urlcolor=blue,
    }
\usepackage{xcolor}

\title{Space-time regularisation}

\begin{document}

\maketitle

\section{Theoretic description}\label{sec:theory}

Let $T\in \mathbb{N}_{>0}$. Let $\Omega\subset\mathbb{R}^d$ be bounded. We wish to reconstruct a series of unknown images $\boldsymbol{I}=(I_0, \ldots, I_T)\in L^2(\Omega)^{T+1}$ using the physical measurements $\boldsymbol{K}=(K_0, \ldots, K_T)$ that are defined by
\begin{equation*}
    \mathcal{F}\circ I_t + \eta_t = K_t \quad\forall 0\le t\le T,
\end{equation*}
where $\mathcal{F}$ denotes the Fourier transform and $\eta_t$ denotes a small error. Thus given $\boldsymbol{K}$ we wish to reconstruct $\boldsymbol{I}$.

We take a velocity field $v\in \mathcal{V}=L^2([0,T], V)$, where $V$ is the closure of $C_c^\infty(\Omega, \mathbb{R}^d)$ w.r.t. the Hölder norm $\|\cdot\|_{C^{1,\alpha}}$ for $0<\alpha\le 1$.

Our deformation of interest is generated by $v$ via the ODE
\begin{align*}
    \frac{\text{d}}{\text{d}t}\varphi(t,x)&= v(t, \varphi(t,x)) & &\forall(t,x)\in [0,T]\times\Omega\\
    \varphi(0,x)&=x & &\forall x\in\Omega.
\end{align*}
We define $\varphi_t=\varphi(t,\cdot)$ and $\varphi_{t_0,t_1}=\varphi_{t_1}\circ\varphi_{t_0}^{-1}$.

Using the following we define our model as:
\begin{equation}\label{energy}
    \mathcal{E}(\boldsymbol{I}, \boldsymbol{\varphi)} = \sum_{t=1}^T \frac12 \|\mathcal{F}\circ I_t - K_t\|_2 + \lambda_1 \mathcal{E}_1(\boldsymbol{I}) + \lambda_2 \mathcal{E}_2 (\boldsymbol{\varphi}) + \lambda_3\mathcal{E}_3 (\boldsymbol{I}, v)
    % \lambda_1 \sum_{t=0}^T \mathcal{R}(I_t) + \lambda_2 \mathcal{E}_v(\boldsymbol{\boldsymbol{I}, \boldsymbol{\varphi}}, v),
\end{equation}
where $\mathcal{E}_1, \mathcal{E}_2, \mathcal{E}_3$ are regularizer terms defined as follows:
\begin{equation*}
    \mathcal{E}_1(\boldsymbol{I}) = \sum_{t=0}^T \mathcal{R}(I_t)
\end{equation*}
where $\mathcal{R}$ is a learned convex ridge regulariser,
\begin{align*}
    \mathcal{E}_2(\boldsymbol{\varphi}) &= \sum_{t=1}^{T-1}\int_\Omega \|\text{D}\varphi_{t,t+1}(x) - \bold{1}\|_2 dx \\
    \mathcal{E}_3(\boldsymbol{I}, v) &= \sum_{t=1}^{T-1}\int_\Omega |\varphi_{t,t+1}\circ I_t - I_{t+1}|^2dx,
\end{align*}

% \begin{equation*}
%     \mathcal{E}_v(\boldsymbol{\boldsymbol{I}, \boldsymbol{\varphi}}, v) = \sum_{t=1}^{T-1}\int_\Omega W(\text{D}\varphi_{t,t+1}(x)) + \lambda_3 \int_t^{t+1}|\text{D}v(\tau, x)|^2d\tau + |\varphi_{t,t+1}\circ I_t - I_{t+1}|^2dx,
% \end{equation*}
% where $W:\mathbb{R}^{d,d}\to \mathbb{R}_{\ge0}$.

\section{Experiment setup involving an inverse problem}

Our goal is to increase reconstruction accuracy by making use of temporal consistency in the measurements.
In many cases, using the additional data directly is not possible since the object moves during the acquisition process, which would introduce a motion blur. 

To make this precise, assume we are given undersampled $k$-space data $K$ (in CMRxRecon, undersampled k-space data for cine MRIs) and the associated forward operator $F$.
We want to solve the inverse problem $F(I) = K$, where $I$ is the unknown image.
The measurement $K$ has a temporal component and shape $(T, H, W)$ as it represents a heartbeat cycle.
In case of CMRxRecon, the ground truth $I_{gt}$ and $K_{gt}$ have resolution $(H_{gt}, W_{gt})$.
The subsampled $K$ has lower resolution, $k\cdot H = H_{gt}$ and $k\cdot W = W_{gt}$ in different tasks, possible values for $k$ are $4, 8, 10$.

Before motion and time dependence are involved, we wish to construct a strong prior at each single time point using no temporal information. We construct this prior using gradient descent and minimizing the optimization objective
\begin{equation}
    \mathcal{E}(\boldsymbol{I}, \boldsymbol{\varphi)} = \sum_{t=1}^T \frac12 \|\mathcal{F}\circ I_t - K_t\|_2 + \lambda_1\mathcal{E}_1(\boldsymbol{I})
\end{equation}
as defined in section \ref{sec:theory} and $\mathcal{R}$ is CRR (with pretrained weights from learned regularisers \href{https://github.com/johertrich/LearnedRegularizers/blob/main/weights/bilevel_CT/CRR_bilevel_JFB_for_CT.pt}{github}, after loading the weights, the scaling is changed to 2).

% \textcolor{cyan}{Option 1:}
% The initial optimization process uses Adam with $\lambda_{1,\text{init}}=10^{-8}$ and learning rate $10^{-2}$, and takes 500 gradient descent step. The reconstruction that achieved the minimal optimization energy is carried forward to the next step. 

% \textcolor{cyan}{Option 2:}
The initial optimization process uses a nonmonotone accelerated proximal gradient method (nmAPG, implementation taken from Learned regularisers \href{https://github.com/johertrich/LearnedRegularizers/blob/main/evaluation/nmAPG.py}{github}). Parameters are: $\lambda_{1, init}=10^{-2}$, max iterations=500, tolerance $10^{-6}$.

The main optimization process has two goals: reconstruction and motion prediction. We have a combined objective w.r.t. which we take gradient descent steps to improve $\boldsymbol{I}$ and $\boldsymbol{v}$ simultaneously.

The combined objective is \ref{energy}.
with the terms described below.
\begin{itemize}
    \item $F$ is the composition of the discrete Fourier transform and a subsampling operator (that subsamples from ($H_{gt}, W_{gt}$ to ($H, W$) using a problem appropriate mask).
    \item $\mathcal{R}$ is the regulariser associated with reconstruction. First choice: convex ridge regulariser (with pretrained weights from learned regularisers \href{https://github.com/johertrich/LearnedRegularizers/blob/main/weights/bilevel_CT/CRR_bilevel_JFB_for_CT.pt}{github}, after loading the weights the scaling is changed to 2).
    \item  $\mathcal{L}_{sim}$ is MSE. Note, that $F$ includes subsampling, thus $F(I_k)$ and $K_k$ have the same shape for each time point $t_k$.
    \item $\boldsymbol{v}$ is the learned velocity field associated with the motion. Output of a Siren network detailed below.
    \item $\varphi$ is the deformation associated with $\boldsymbol{v}$. It is the solution of the ODE $\frac{d\varphi}{dt}(t,x) = \boldsymbol{v}(\varphi(t,x),t)$ as described in \ref{sec:theory}.
    % \item $\mathcal{L}_{\text{motion reg}}$ is the regulariser of $\boldsymbol{v}$, described more precisely below.
\end{itemize}

The non-stationary velocity field $\boldsymbol{v}(x,t)$ is generated by a Siren implicit neural representation network at each time point, so essentially an ensemble of Siren networks. Network $k$ produces the velocity for the time interval $[k,k+1)$. Then we solve the ODE. Backpropagation through the ODE is done using the adjoint method. We use explicit Euler to solve the ODE.
% The motion regulariser $\mathcal{L}_{\text{motion reg}}$ consists of two parts: velocity gradient and negative deformation Jacobian determinant regularisation
% \begin{itemize}
%     \item Velocity gradient regulariser: $\mathcal{L}_{grd}=\frac{1}{N\cdot T}\sum_{t, x} \| \nabla\boldsymbol{v_\theta}(x,t)\|_2^2$
%     \item Regulariser that penalizes the negative determinants of the Jacobian of the deformation $\varphi$: $\mathcal{L}_{detJ} = \frac{1}{N\cdot T} \sum_{t,x} \log( \text{ReLU}(-\det \boldsymbol{J}_{\varphi_t}(x) + 1))$
% \end{itemize}

\subsection{Hyperparameters}

Important hyperparameters to choose:
\begin{itemize}
    \item Multiple steps in the ODE solver between each consecutive time point $t_i$ and $t_{i+1}$. First try will be 10. Thus $10(T-1)$ steps in total.
    \item Exact choice of network for the velocity. Per Sebastian's recommendation, I will try an ensemble of Siren networks the following way: between each consecutive time points the velocity will be stationary, represented by a purely spatial INR. But in different time intervals the networks are different so overall the velocity is non-stationary.
    \item Initialisation of the Siren networks. The velocity is mostly 0 thus a small initial velocity is probably beneficial.
\end{itemize}


\section{Datasets}

\subsection{Time dependent datasets/indirect observations}

\begin{itemize}
    \item \href{https://www.nature.com/articles/s41597-024-03525-4}{CMRxRecon}: Cine MRI scans. Full k-space is given but masks are available to simulate an undersampled k-space.

    \item \href{https://fips.fi/open-datasets/x-ray-tomographic-datasets/stempo-dynamic-x-ray-tomography-phantom-dataset/}{STEMPO}: real X-ray tomography data measured with a motorized device. Dataset used in \href{https://arxiv.org/pdf/2405.06337}{Regularization with optimal space-time priors}.

    \item \href{https://aimi.stanford.edu/datasets/sinoct}{SinoCT} Dataset of 9000 head CTs and corresponding (simulated) sinograms.

    \item \href{https://med.emory.edu/departments/radiation-oncology/research-laboratories/deformable-image-registration/index.html}{DIR-LAB} 4DCT thoracic images and exhale-inhale CT image pairs

    \item \href{https://image-x.sydney.edu.au/spare-challenge/}{SPARE challange} 4D cone-beam CTs

    \item \href{https://www.cancerimagingarchive.net/collection/4d-lung/}{4D-LUNG}: 4D fan beam and cone beam CTs.

    \item \href{https://www.ub.edu/mnms/}{Multi-Centre, Multi-Vendor \& Multi-Disease Cardiac Image Segmentation Challenge (M\&Ms) }: 4D short-axis cardiac MRIs. Annotations are provided for 150 studies for only two cardiac time frames, end-diastole (ED) and end-systole (ES).

    \item \href{https://www.creatis.insa-lyon.fr/Challenge/acdc/databases.html}{ACDC} 2D cine-MRIs, each case contains scans at two time frames, one at end-diastole (ED) and one at end-systole (ES). All scans come with segmentation of the right and left ventricle cavities, as well as the myocardium.
\end{itemize}

\subsection{Purely spatial datasets}

\begin{itemize}
    \item Brain MRIs: \href{http://braintumorsegmentation.org/}{BraTS challange}, \href{https://sites.wustl.edu/oasisbrains/home/oasis-1/}{OASIS-1}, \href{https://www.nitrc.org/projects/candi_share/}{CANDI}, \href{https://adni.loni.usc.edu/data-samples/adni-data/neuroimaging/mri/}{ADNI}, \href{https://mindboggle.info/data}{Mindboggle-101}, \href{https://www.loni.usc.edu/research/atlases}{LPBA40}

    \item PET-CT: \href{https://autopet-iii.grand-challenge.org/task/}{autoPET}
\end{itemize}

\section{Papers of interest}

\subsection{Neural fields}

\begin{itemize}
    \item \href{https://arxiv.org/abs/2305.06822}{Implicit Neural Networks with Fourier-Feature Inputs for Free-breathing Cardiac MRI Reconstruction}
    \item \href{https://arxiv.org/abs/2307.14363}{Unsupervised reconstruction of accelerated cardiac cine MRI using Neural Fields}
    \item \href{https://openaccess.thecvf.com/content/ICCV2021/html/Reed_Dynamic_CT_Reconstruction_From_Limited_Views_With_Implicit_Neural_Representations_ICCV_2021_paper.html}{Dynamic CT Reconstruction From Limited Views With Implicit Neural Representations and Parametric Motion Fields}
    \item \href{https://arxiv.org/pdf/2403.03860}{ProxNF: Neural Field Proximal Training for High-Resolution 4D Dynamic Image Reconstruction}
\end{itemize}

\subsection{General ODE based diffeomorphic image registration}

\begin{itemize}

\item \href{https://arxiv.org/abs/2305.14673}{ORRN: An ODE-Based Recursive Registration Network for Deformable Respiratory Motion Estimation With Lung 4DCT Images}

2021, 11 citations, \href{https://github.com/Lancial/orrn_public}{github}, DS: CREATIS, DIR-Lab

Needs training set

The novelty is in the architecture predicting $v$. Can handle time-series of 3D images. Allows for single and multi resolution architecture. Claims that the multi-resolution case is superior but only downsamples once or twice, and admits that this results in increased computational costs.

\item \href{https://arxiv.org/abs/2106.08493}{Multi-scale Neural ODEs for 3D Medical Image Registration}

2021, 30 citations, no code available, DS: BraTS

Exact architecture is very poorly described. Multi-scale ODE, as time goes on, the feature maps are downsampled. They work on multimodal registration so the similarity is calculated in a modality independent space after feature extraction by a pretrained net.

\item \href{https://proceedings.mlr.press/v172/wolterink22a.html}{Implicit Neural Representations for Deformable Image Registration}

2022, 99 citations, \href{https://github.com/MIAGroupUT/IDIR}{github}, DS: DIR-Lab

Uses a lightweight MLP to predict the transformation of a pixel, essentially for an input coordinate it outputs another one. Can be analytically differentiated if not using ReLU. Is worth having a look at.

\item \href{https://arxiv.org/pdf/2102.07951}{ResNet-LDDMM: Advancing the LDDMM
Framework Using Deep Residual Networks}

2022, 39 citations, no code available, DS: SHREC

LDDMM inspired registration method, it uses ResNet blocks to solve the flow equation, otherwise no changes are made. The method is faster than LDDMM especially for high resolution and the authors claim that it has an advantage comapred to LDDMM because there is no need to find a suitable scale.

\item \href{https://cnrs.hal.science/hal-03971473/file/WBIR_2022_WeightedMetamorphosis.pdf}{Weighted Metamorphosis for registration of
images with different topologies}

2022, 11 citations, \href{https://github.com/antonfrancois/Demeter_metamorphosis}{github}, DS: BraTS, synthetic

Uses a time-dependent spatial weight on the intensity. This helps disentangle geometric and intensity changes. Uses prior knowledge to construct such weights, i.e. it is zero outside of the segmentation of a tumor.

\item \href{https://www.sciencedirect.com/science/article/pii/S1361841523001779}{R2Net: Efficient and flexible diffeomorphic image registration using Lipschitz continuous residual networks}

2023, 17 citations, \href{https://github.com/ankitajoshi15/R2Net}{github}, DS: ACDC, EMPIRE10, OASIS, IBSR18

R2Net uses a UNet to produce the initial velocity field ($v_0$). A series of consecutive Lipschitz continuous ResNet blocks model the solution of the flow equation (the output of the $k^{th}$ block is the velocity field at time $k$). The ResNet blocks can either share their parameters (stationary velocity field) or not.

\item \href{https://openaccess.thecvf.com/content/WACV2023/html/Han_Diffeomorphic_Image_Registration_With_Neural_Velocity_Field_WACV_2023_paper.html}{Diffeomorphic Image Registration With Neural Velocity Field}

2023, 28 citations, no code available

\item \href{https://arxiv.org/abs/2303.04849}{MetaMorph: Learning Metamorphic Image Transformation with Appearance Changes}

2023, 10 citations, \href{https://github.com/jw4hv/metamorph/blob/Initial/README.md}{github}

\item \href{https://openaccess.thecvf.com/content/CVPR2022/papers/Wu_NODEO_A_Neural_Ordinary_Differential_Equation_Based_Optimization_Framework_for_CVPR_2022_paper.pdf}{NODEO: A Neural Ordinary Differential Equation Based Optimization
Framework for Deformable Image Registration}

2024, 40 citations, \href{https://github.com/yifannnwu/NODEO-DIR}{github}

\item \href{https://www.sciencedirect.com/science/article/pii/S1361841524001749}{Medical image registration via neural fields}

2024, 23 citation, no code available

\item \href{https://arxiv.org/abs/2401.17493}{CLAIRE: Scalable GPU-Accelerated Algorithms for Diffeomorphic Image Registration in 3D} 

2024, \href{https://github.com/andreasmang/claire}{github}

\item \href{https://arxiv.org/abs/2405.18684}{Learning Diffeomorphism for Image Registration with Time-Continuous Networks using Semigroup Regularization}

2024, 1 citation, \href{https://github.com/mattkia/SGDIR}{github}

\item \href{https://arxiv.org/abs/2407.13241}{NODER: Image Sequence Regression Based on Neural Ordinary Differential Equations}

2024, 1 citation, \href{https://github.com/ZedKing12138/NODER-pytorch}{github}

\item \href{https://arxiv.org/abs/2405.15385}{CPT-Interp: Continuous sPatial and Temporal Motion Modeling for 4D Medical Image Interpolation}

2024, no code available

\item \href{https://arxiv.org/pdf/2412.17982}{Unsupervised learning of spatially varying regularization for diffeomorphic image
registration}

2024, \href{https://github.com/junyuchen245/Spatially-Varying-Regularization-ImgReg}{github}

\item \href{https://www.sciencedirect.com/science/article/pii/S1877750324003004}{PDE-LDDMM meets NODEs: Introducing neural ordinary differential equation solvers in PDE-constrained Large Deformation Diffeomorphic Metric Mapping}

2025, no code available

NODEO-LDDMM: changes NODEO by including an extra regularizer term inspired by LDDMM $\int_0^1 \| v_t^{\theta}(\phi _{t,0})\|_V^2dt$. $v_t^\theta$ is the solution to the transport equation estimated by the NODE.

\item \href{https://arxiv.org/abs/2503.13756}{Fast alignment of heterogeneous images in sliced Wasserstein distance}

2025

\end{itemize}

\section{Old baselines}

\begin{itemize}
    \item \href{https://www.sciencedirect.com/science/article/pii/S1361841507000606}{SyN} (2008), \href{https://docs.dipy.org/stable/examples_built/registration/syn_registration_2d.html#footcite-avants2008}{implementation}

    \item  \href{https://arxiv.org/pdf/1809.05231}{VoxelMorph} (2017), \href{https://github.com/voxelmorph/voxelmorph/tree/dev}{github}, (or \href{https://arxiv.org/pdf/2101.01035}{HyperMorph} or \href{https://github.com/junyuchen245/TransMorph_Transformer_for_Medical_Image_Registration}{TransMorph)}
\end{itemize}

\section{Survey}

\href{https://www.sciencedirect.com/science/article/pii/S1361841524003104?ref=pdf_download&fr=RR-2&rr=8fdb47670bc2d390}{A survey on deep learning in medical image registration: New technologies, uncertainty, evaluation metrics, and beyond}

\href{https://github.com/JHU-MedImage-Reg/MedImgReg_Survey}{github}



\end{document}
